\section{Information Growth Across Channel Regimes}
\label{sec:pivotal-middle}

\subsection{Intermediate-Correlation Growth}

We first prove the required growth inequality on the intermediate range
\(31/20\le\ell=\atanh\rho\le7/2\). In this section write \(r=\rho\).
The idea is to keep the coordinates together. Logarithmic Sobolev gives
an entropy cost for their pivotal probabilities; the log-sum inequality
then combines that cost with the edge budget. We will not choose a
comparison Fourier degree or divide coordinates by influence.
The underlying entropy tool is Gross's cube logarithmic Sobolev
inequality~\cite{Gross,ODonnell}. Its use here is specific: the positive
kernel in the next calculation converts pivotal-set entropy into the
same spectral expression produced by the centered angle comparison.
This supplies the correction left after the entropy-flow edge
bound~\cite{CGN}, without a conjectured coupling potential.


\begin{proposition}[Pivotal entropy forces faster growth]
\label{prop:middle-growth}
Let \(f\) be monotone Boolean and
\(31/20\le\atanh\rho\le7/2\). Then
\[
 I_f(\rho)>\phi(\rho)\quad\Longrightarrow\quad
 I_f'(\rho)>\phi'(\rho).
\]
The mean of \(f\) is unrestricted.
\end{proposition}

The proof compares the energy demanded by the information premise with
the growth available along the cube edges. It first obtains an
angle-energy requirement, expresses part of that requirement as pivotal
entropy, and then combines the entropy with the edge correction through
the log-sum inequality.

Put \(\Theta=\arcsin u\) and \(J_{\rm ang}=\E[\Theta\Lapl\Theta]\).
The dictator has posterior values \(\pm r\). Define its angle and the
slope of the line through those two angle values by
\[
 \theta=\arcsin r,\qquad \lambda=\theta/r.
\]
Two further coefficients come from the tangent inequality below:
\[
 \alpha=\frac{\theta+\sinh\ell}{\ell},\qquad
 \gamma=\lambda\alpha,\qquad
 \eta=\frac{\theta^2\gamma}{\gamma+1}.
\]
These are channel functions, not parameters to optimize.
The same \(\gamma\) will serve throughout the growth proof, including
the high-correlation extension. Only the way we estimate the pivotal
cost changes with the channel.

\subsubsection{Angle energy required by an information lead}

The following scalar inequality is proved in
Appendix~\ref{app:centered-support}:
\begin{equation}
 \begin{split}
 \frac{\alpha}{2}(\arcsin x-\lambda x)^2
 &\le x\arcsin x-r\theta\\
 &\quad-\alpha[\phi(x)-\phi(r)],qquad -1\le x\le1.
 \end{split}
 \label{eq:centered-support}
\end{equation}
Both sides vanish at \(x=\pm r\). It bounds the error of the
dictator's angle line by an entropy difference, without losing
equality at the reference posterior.
The comparison stays within Shannon entropy. Although the angle and
the Hellinger program~\cite{Hellinger,CN} both involve the geometry of
binary probabilities, we do not infer this inequality from a
Hellinger-optimality conjecture; its scalar proof is given in full.

Average \eqref{eq:centered-support} with \(x=u\), multiply by
\(2\lambda\), and complete the square in the positive operator
\(\Lapl+\gamma\). This gives
\begin{equation}
 \begin{split}
 J_{\rm ang}\ge{}&2\theta^2+2\gamma[\phi(m)+\Delta_f(r)]\\
 &+\lambda^2\E\left[u\left\{\gamma-
       \frac{(\gamma+1)^2}{\Lapl+\gamma}\right\}u\right].
 \end{split}
 \label{eq:centered-resolvent}
\end{equation}
Here division by \(\Lapl+\gamma\) means multiplication of each
degree-\(k\) Fourier coefficient by \(1/(k+\gamma)\).
For completeness, the discarded nonnegative square is
\begin{multline*}
 \E\bigl[
 \bigl(\Theta-\lambda(\gamma+1)(\Lapl+\gamma)^{-1}u\bigr)
 (\Lapl+\gamma)\\
 {}cdot\bigl(\Theta-\lambda(\gamma+1)
 (\Lapl+\gamma)^{-1}u\bigr)\bigr].
\end{multline*}
The constant Fourier mode is included; this matters when \(m\ne0\).
For the dictator, \(u=ry_1\) and \(\Theta=\lambda u\) are both
degree one. The operator \((\gamma+1)(\Lapl+\gamma)^{-1}\)
therefore fixes \(u\), so the discarded square is zero. In
\eqref{eq:centered-resolvent} the spectral term is \(-\theta^2\);
the result is exactly \(J_{\rm ang}=\theta^2\).

\subsubsection{Contribution of noisy pivotal sets}

Recall the pivotal indicator \(q_i^f\), with mean \(a_i\).
Omit coordinates with \(a_i=0\); every logarithm below then has a
positive argument.
Let \(t\in[0,r^2]\) denote a squared noise correlation and put
\[
 w_i(t)=T_{\sqrt t}q_i^f,\qquad s_i(t)=\E w_i(t)^2.
\]
Noise here acts on the other \(n-1\) coordinates.
The value \(w_i(t)\) is the conditional probability that coordinate
\(i\) is pivotal. Its mean remains \(a_i\); its second moment
\(s_i(t)\) records how unevenly that probability is distributed.
Complete averaging gives \(s_i(0)=a_i^2\), whereas the unsmoothed
indicator has second moment \(a_i\).
Define an average of these second moments:
\[
 z_i=\frac{\gamma+1}{\gamma r^2}
       \int_0^{r^2}[1-(t/r^2)^\gamma]s_i(t)\,dt,
 \qquad Z=\sum_i z_i.
\]
The weight in this integral is nonnegative and has total mass one.
Thus \(a_i^2\le z_i\le a_i\): noise preserves the mean and cannot
increase the second moment of an indicator.

Write \(V_0=1-m^2\) for the Boolean variance and
\(U=\operatorname{Var}(u)/r^2\) for the normalized noisy variance.
Fourier expansion of the integral gives
\[
 Z=(\gamma+1)\sum_{k\ge1}
       \frac{r^{2(k-1)}}{\gamma+k}
       \sum_{|S|=k}\widehat f(S)^2,
 \qquad 0<Z\le U\le V_0.
\]
The last inequality is Boolean Parseval followed by noise contraction.
For a dictator, \(w_1(t)=s_1(t)=z_1=1\), all other
coordinate terms vanish, and \(Z=U=V_0=1\). Thus \(V_0-Z\)
measures a variance reserve that is absent in the equality example.
The exact integral identity
\begin{align*}
 &\frac{\gamma+1}{\gamma}
 \int_0^{r^2}[1-(t/r^2)^\gamma]
       \sum_i\E[w_i(t)\Lapl w_i(t)]\,dt\\
 &\hspace{7em}=(\gamma+1)r^2(U-Z)
\end{align*}
is checked one Fourier degree at a time: a degree-\(k\) coefficient
appears in \(k\) pivotal derivatives, each of degree \(k-1\).

Cube logarithmic Sobolev, proved in Appendix~\ref{app:scalar}, gives
\[
 \E[w_i\Lapl w_i]\ge
       \frac12 s_i\log\frac{s_i}{a_i^2}.
\]
For the auxiliary entropy bound, average under the probability
density \(w_i^2/s_i\). Jensen applied to \(\log(1/w_i)\) gives
\(\E_*\log w_i\ge\log(s_i/a_i)\), hence
\(\operatorname{Ent}(w_i^2)\ge s_i\log(s_i/a_i^2)\).
Here \(\operatorname{Ent}(v)=\E(v\log v)-(\E v)\log(\E v)\);
zero-weight points are omitted. Convexity of \(s\log(s/a_i^2)\)
and the averaging integral therefore give a bound on the
\emph{pivotal entropy cost}, defined by
\[
 E_{\rm piv}:=\sum_i z_i\log\frac{z_i}{a_i^2}.
\]
Each summand is nonnegative since \(z_i\ge a_i^2\); a noisy pivotal
probability that is constant has zero cost. In this notation,
\begin{equation}
 (\gamma+1)(U-Z)\ge\frac12 E_{\rm piv}.
 \label{eq:pivotal-entropy}
\end{equation}

Rewrite \eqref{eq:centered-resolvent} in terms of \(U,Z,V_0\), including
its constant mode, and then apply \eqref{eq:pivotal-entropy}.
To keep the bias contribution visible, define
\[
 M(m):=2\gamma\phi(m)-[\lambda^2(2-r^2)+\lambda/\alpha]m^2.
\]
The required angle energy is now
\begin{align}
 J_{\rm ang}
 &\ge\theta^2[1+V_0+\gamma U-(\gamma+1)Z]
          +M(m)+2\gamma\Delta_f(r)\notag\\
 &\ge\theta^2(1+V_0-Z)+\frac{\eta}{2}E_{\rm piv}
          +M(m)+2\gamma\Delta_f(r).                 \label{eq:pivotal-baseline}
\end{align}
Read the last line term by term: the dictator's energy \(\theta^2\),
the variance reserve \(\theta^2(V_0-Z)\), the pivotal entropy cost,
the bias correction, and the credit \(2\gamma\Delta_f(r)\) for extra
information. The last credit is strictly positive under a lead.
For the dictator, \(E_{\rm piv}=1\log1=0\), and every term except
\(\theta^2\) vanishes.

\subsubsection{Log-sum closure of the edge budget}

Suppose the information lead \(\Delta=\Delta_f(r)\) is positive,
but its growth is no faster than the dictator's: \(A(u)\le r\ell\).
We derive a contradiction.
Monotonicity gives \(a_i\le1-|m|\le V_0\). Also
\[
 \E h(u)=h(r)-\phi(m)-\Delta\le V_0h(r),
\]
since \(\phi(m)\ge m^2/2\) and \(h(r)<1/2\).
For the latter, \(r>9/10\) and \(L<7/10\) give
\(h(r)\le L-r^2/2<1/2\).
Apply Lemma~\ref{lem:edges} with
\(H_{\rm budget}=V_0h(r)\). Its heights satisfy
\[
 F(v_i)=\frac{V_0h(r)}{r a_i}=\frac{V_0F(\ell)}{a_i},
 \qquad 0<v_i\le\ell,
\]
where the upper bound follows from \(a_i\le V_0\) and the decrease
of \(F\). With \(H(v)=v-c(v)\), the lemma gives
\begin{equation}
 \sum_i a_i v_i\le A(u)/r\le\ell,\qquad
 A(u)\ge J_{\rm ang}+r\sum_i a_iH(v_i).
 \label{eq:adjusted-edge}
\end{equation}
These heights are comparison parameters; they are not additional
properties assumed of the field.

Set \(\beta=2r/\eta\). The channel estimates proved in
Lemma~\ref{lem:middle-channel-checks} are
\begin{align}
 \log\frac{vF(v)}{\ell F(\ell)}
       &\ge\beta[H(\ell)-H(v)] &&(0<v\le\ell), \label{eq:middle-profile}\\
 \theta^2-rH(\ell)-\eta/2&\ge0,                \label{eq:middle-energy}\\
 M(m)-rH(\ell)m^2&\ge0.                       \label{eq:middle-mean}
\end{align}
Their role is simple: the first converts the edge budget to a
normalization for log-sum; the other two leave nonnegative reserves.

For that normalization, put
\(b_i=a_i^2e^{-\beta H(v_i)}\) and \(B_{\rm piv}=\sum_i b_i\).
This choice puts the edge correction inside the entropy logarithm:
\(\log(z_i/b_i)=\log(z_i/a_i^2)+\beta H(v_i)\).
The weights need not sum to one; log-sum keeps track of their total.
Equation~\eqref{eq:middle-profile} gives
\[
 \frac{b_i}{a_iv_i}
 \le\frac{V_0}{\ell}e^{-\beta H(\ell)},\qquad
 B_{\rm piv}\le V_0e^{-\beta H(\ell)}.
\]
The second inequality uses the first bound in
\eqref{eq:adjusted-edge}: a rule growing no faster than the dictator
has only \(\ell\) units of edge weight to allocate.
Since \(z_i\le a_i\), \(H\ge0\), and
\(\eta\beta/2=r\), log-sum now yields
\begin{align*}
 \frac{\eta}{2}E_{\rm piv}+r\sum_i a_iH(v_i)
 &\ge\frac{\eta}{2}\sum_i z_i\log\frac{z_i}{b_i}\\
 &\ge\frac{\eta}{2}Z\log\frac Z{B_{\rm piv}}\\
 &\ge\frac{\eta}{2}Z\log\frac Z{V_0}+rZH(\ell).
\end{align*}
The first line joins the two costs by the logarithmic identity just
given. The second is log-sum. The third uses the bound on total weight.
This is the connection between the information requirement and the
available growth: both involve the same noisy pivotal sets.

Finally insert this bound into \eqref{eq:pivotal-baseline} and
\eqref{eq:adjusted-edge}. Subtract
\(r\ell=\theta^2+rH(\ell)\), and use
\(Z\log(Z/V_0)\ge Z-V_0\), the tangent inequality for \(t\log t\):
\begin{equation}
 \begin{aligned}
 A(u)-r\ell\ge{}&
 \underbrace{[\theta^2-rH(\ell)-\eta/2](V_0-Z)}_{
                    \text{nonnegative variance reserve}}\\
 &+\underbrace{[M(m)-rH(\ell)m^2]}_{\text{nonnegative bias reserve}}
  +\underbrace{2\gamma\Delta}_{\text{positive lead credit}}>0.
 \end{aligned}
 \label{eq:pivotal-growth-account}
\end{equation}
This contradicts \(A(u)\le r\ell\). We have proved
\[
 \Delta_f(r)>0\ \Longrightarrow\ A(u)>r\ell
 \qquad(31/20\le\ell\le7/2).
\]
At a dictator, \(a_1=z_1=V_0=1\) and \(v_1=\ell\).
Log-sum has one term and is exact. The variance reserve, bias reserve,
and lead credit are all zero. What remains is the exact account
\[
 \underbrace{r\ell}_{\text{growth}}
   =\underbrace{\theta^2}_{\text{angle energy}}
    +\underbrace{rH(\ell)}_{\text{edge correction}}.
\]
An information lead would add the strictly positive last term in
\eqref{eq:pivotal-growth-account}; the other two terms cannot cancel it.
